\section{Introduction}
\label{sec:introduction}

Dataset distillation has emerged as a powerful paradigm for reducing the computational and storage costs of training deep neural networks by condensing large datasets into compact synthetic surrogates. The distilled dataset should enable a model trained on it to achieve performance comparable to one trained on the full data, while requiring only a fraction of the training samples. Early distillation methods formulated this as a bi-level optimization problem, alternating between updating synthetic data and training a surrogate model via gradient matching~\cite{zhao2023improved}, trajectory matching~\cite{cazenavette2022dataset}, or distribution matching~\cite{liu2023dream}. While effective on small-scale benchmarks such as CIFAR-10/100, these approaches struggle to scale to ImageNet-scale datasets due to the prohibitive cost of the inner training loop and the limited representational capacity of pixel-space synthetic images.

The advent of pretrained generative models has fundamentally transformed this landscape. Methods such as SRe$^2$L~\cite{yin2023squeeze} decoupled the bi-level optimization by first training a generative model on the real data and then sampling unlimited synthetic images, enabling distillation at ImageNet scale. More recently, MGD$^3$~\cite{chansantiago2025modeguided} demonstrated that a pretrained class-conditional diffusion model---specifically, a Diffusion Transformer (DiT)~\cite{peebles2022scalable} trained on ImageNet-1K---can serve as a direct, training-free generator for dataset distillation. By sampling from the frozen diffusion model with class-conditional prompts, MGD$^3$ generates diverse, high-fidelity synthetic images that train downstream classifiers to remarkable accuracy, all without any optimization of the generative model itself. This represents a paradigm shift: the distilled dataset is not optimized but rather \emph{sampled} from a pretrained prior.

Despite its effectiveness, MGD$^3$ and its contemporaries apply a single, globally fixed classifier-free guidance (CFG)~\cite{ho2022classifierfree} strength to all classes and all denoising timesteps. Classifier-free guidance interpolates between the conditional and unconditional score functions via a scalar $\lambda$, and the choice of $\lambda$ critically governs the quality--diversity trade-off of the generated samples. A low $\lambda$ produces diverse but low-fidelity images that may fail to capture class-discriminative features, while an excessive $\lambda$ yields high-fidelity but mode-collapsed samples that lack the intra-class variability necessary for robust classifier training. Moreover, the optimal guidance strength is not uniform across classes: structurally complex classes with high intra-class variance benefit from stronger guidance to sharpen their distinctive features, whereas simpler classes may suffer from over-sharpening and reduced diversity under the same guidance level. This per-class heterogeneity is entirely ignored when a single $\lambda$ is applied globally.

We address these limitations with \ags{} (Adaptive Guidance Scheduling for Diffusion Dataset Distillation), a training-free framework whose core module, \cags{}, replaces the globally fixed guidance strength with per-class adaptive guidance:

\begin{itemize}
    \item \textbf{\cags{} (Class-Adaptive Guidance Strength):} Assigns a per-class guidance magnitude based on a multi-factor complexity metric that quantifies mode count, cluster entropy, intra-class variance, and inter-class separability. Through systematic factor ablation (Section~\ref{sec:factor-ablation}), we find that \emph{per-class adaptation}---not merely the presence of guidance---drives the improvement: cluster entropy is the most effective single complexity factor, while inter-class separability---which produces near-constant guidance across classes---is the least effective. We provide a theoretical explanation (Proposition~\ref{prop:entropy}): under CFG, high-entropy (balanced multi-modal) classes are more sensitive to the guidance scale because raising the conditional density to the power $(1+\lambda)$ disproportionately collapses their smaller modes, making them more likely to benefit from adaptive rather than fixed guidance. Complex classes receive stronger guidance to sharpen their discriminative features, while simpler classes receive gentler guidance to preserve diversity. \cags{} yields substantial improvements on the 100-class subset ($+7.6\%$ to $+14.8\%$ relative) and at 10-class IPC=50 ($+9.0\%$), with gains scaling on deeper architectures ($+17.0\%$ on ResNet-18).
    \item \textbf{Exploratory modules (not recommended):} We additionally explore two complementary modules---\iast{} (IPC-Adaptive Guidance Range) and \tags{} (Timestep-Adaptive Guidance Schedule)---and report their negative ablation results as boundary conditions. \iast{}'s log-scaling is counterproductive at low IPC, and \tags{} does not improve over \cags{}-alone, confirming that the class axis---not the temporal or budget axis---is the primary source of suboptimality in fixed-guidance distillation.
\end{itemize}

\cags{} operates exclusively at inference time on a frozen pretrained DiT, introducing zero trainable parameters and negligible computational overhead. This distinguishes \ags{} from training-based approaches such as Learnability-Guided Diffusion (LGD)~\cite{chansantiago2026learnabilityguided}, which requires iterative, curriculum-based sample selection over multiple training rounds. While LGD achieves strong results by learning to select the most learnable samples, it sacrifices the simplicity and speed of single-pass generation. \ags{} preserves the training-free, single-pass paradigm of MGD$^3$ while substantially improving the quality of the generated dataset through adaptive guidance.

We validate \ags{} through extensive experiments on ImageNet-100 (both 10-class and 100-class subsets), ImageNette, ImageWoof, semantically coherent subsets of ImageNet-100, and a 200-class ImageNet subset, across ConvNet-6, ResNet-18, and ResNetAP-10 architectures, at IPC settings of 1, 10, and 50, all under the MGD$^3$-matched evaluation protocol with best-accuracy tracking. Our results reveal three key findings. First, \emph{\cags{} provides substantial and consistent gains on the 100-class subset}: $+7.6\%$ relative at IPC=10 and $+14.8\%$ at IPC=50, with gains scaling on deeper architectures ($+17.0\%$ on ResNet-18, $+13.4\%$ on ResNetAP-10). Fixed guidance \emph{hurts} on this subset ($-1.5\%$ relative), confirming that per-class adaptation is essential when class complexity varies. Scaling to 200 classes amplifies the \cags{} advantage from $+2.50\%$ to $+3.68\%$ absolute, confirming that per-class adaptation becomes increasingly beneficial as the complexity distribution widens. A fair comparison with D4M under the same DiT generator reveals that real-image initialization is sensitive to the denoising start timestep ($t_{\text{start}}$): at $t_{\text{start}}{=}25$ it underperforms unguided ($20.21\%$ vs.\ $21.63\%$), while at $t_{\text{start}}{=}40$ it approaches \cags{} ($23.07\%$ vs.\ $23.27\%$) but requires real images and hyperparameter tuning, confirming that \cags{}'s training-free approach is more principled. Second, \emph{the advantage of multi-factor \cags{} over entropy-only guidance is scale-dependent}: on small-scale datasets ($\leq$10 classes), entropy-only guidance outperforms the full multi-factor metric by $1.4$--$1.6\%$, while on medium-to-large-scale datasets ($\geq$44 classes), the multi-factor \cags{} metric wins by $0.65$--$0.94\%$. A semantic subset analysis---partitioning ImageNet-100 into 44 animal and 56 non-animal classes---confirms that the multi-factor advantage is driven by class count, not semantic composition. Third, \emph{per-class adaptation is the key mechanism, not merely the presence of guidance}: a systematic ablation over all four complexity factors reveals that cluster entropy is the most effective single factor, while inter-class separability---which produces near-constant guidance (std $0.006$)---is the least effective. We provide a theoretical explanation (Proposition~\ref{prop:entropy}) linking cluster entropy to CFG's mode-collapse behavior: high-entropy (balanced multi-modal) classes are more sensitive to the guidance scale because raising the conditional density to the power $(1+\lambda)$ disproportionately collapses their smaller modes, making them more likely to benefit from adaptive rather than fixed guidance. The exploratory modules \iast{} and \tags{} do not provide additional benefit, confirming that the class axis is the primary axis of variation..
